%------------------------------------------------------------------------------%

\usepackage{integracao_e_series}
\usepackage{ulem}

\title{Integração por Substituição Trigonométrica}
\subtitle{3 -- Substituição pela secante}

%------------------------------------------------------------------------------%
\begin{document}

\addtitle

%------------------------------------------------------------------------------%
\section{Integração por Substituição Trigonométrica}

%------------------------------------------------------------------------------%
\begin{frame}{Substituições Trigonométricas}
\addtext{16}{0.35,2.2}{
\renewcommand*{\arraystretch}{2.3}
\begin{tabular}{c<{\;}>{\;}c<{\;}>{\;}c<{\;}>{\;}c} \hline
  Expressão
  & Substituição
  & Intervalo
  & Identidade \\ \hline
  \(\sqrt{a^2-x^2}\)
  & \(x=a\sin(\theta)\)
  & \(-\frac{\pi}{2} \leq \theta \leq \frac{\pi}{2}\)
  & \(1-\sin^2(\theta) = \cos^2(\theta)\) \\
  \(\sqrt{a^2+x^2}\)
  & \(x=a\tan(\theta)\)
  & \(-\frac{\pi}{2}< \theta < \frac{\pi}{2} \)
  & \(1+\tan^2(\theta)= \sec^2(\theta)\) \\
  {\color{blue}\(\sqrt{x^2-a^2}\)}
  & {\color{blue}\(x=a\sec(\theta)\)}
  & \(0 \leq \theta < \frac{\pi}{2} \) ou \(\pi \leq \theta < \frac{3\pi}{2} \)
  & \(\sec^2(\theta)-1= \tan^2(\theta)\) \\ \hline
\end{tabular}}
\end{frame}

% 1
%------------------------------------------------------------------------------%
%------------------------------------------------------------------------------%
\addtocounter{exemplo}{1}
\section{Exemplo \theexemplo}

%------------------------------------------------------------------------------%
\begin{frame}{Exemplo \theexemplo}
  Encontre \(\int\!\!\frac{dx}{\sqrt{x^2-a^2}}\) com $a>0$

  \pause
  \vspace{0.5\baselineskip}
  Substituição
  \[
    x = a\sec(\theta)
    \qquad
    dx = a \sec(\theta)\tan(\theta) d\theta
  \]
  onde \(0 < \theta < \frac{\pi}{2}\) ou \(\pi < \theta  < \frac{3\pi}{2}\)

  \pause
  \vspace{0.5\baselineskip}
  \[
    F = \int \frac{ a \sec(\theta) \tan(\theta) }
                  { \sqrt{a^2 \sec^2(\theta) - a^2 }} d\theta
  \]
\end{frame}

%------------------------------------------------------------------------------%
\begin{frame}{Exemplo \theexemplo}
\begin{columns}
\column{0.5\linewidth}
  \begin{align*}
    \onslide<+->{F & = \int \frac{dx}{\sqrt{x^2-a^2}}                                       \\[4mm]}
    \onslide<+->{  & = \int  \frac{ a \sec(\theta) \tan(\theta)}{\sqrt{a^2 \sec^2(\theta)-a^2}}  d\theta \\[4mm]}
    \onslide<+->{  & = \int \frac{a \sec(\theta) \tan(\theta)}{a\sqrt{\sec^2(\theta)-1}}  d\theta}
  \end{align*}
\column{0.5\linewidth}
  \onslide<+->{Usando \(\sec^2(\theta)-1= \tan^2(\theta)\)}
  \begin{align*}
    \onslide<+->{F & = \int \frac{\sec(\theta) \tan(\theta)}{\sqrt{\tan^2(\theta)}} d\theta \\[4mm]}
    \onslide<+->{  & = \int \frac{\sec(\theta)\tan(\theta)}{\tan(\theta)} d\theta           \\[4mm]}
    \onslide<+->{  & = \int \sec(\theta) d\theta                                            }
  \end{align*}
\end{columns}
\end{frame}

%------------------------------------------------------------------------------%
\begin{frame}{Exemplo \theexemplo}
  \begin{align*}
    \onslide<+->{F & = \int \sec(\theta) d\theta \\[2mm]}
    \onslide<+->{  & = \int \sec(\theta)
                          \;\frac{\sec(\theta)+\tan(\theta)}
                                 {\sec(\theta)+\tan(\theta)}\; d\theta \\[2mm]}
    \onslide<+->{  & = \int \frac{\sec^2(\theta)+ \sec(\theta)\tan(\theta)}
                                 {\sec(\theta)+\tan(\theta)} d\theta}
  \end{align*}

  \onslide<+->{Substituição
  \(
    u = \sec(\theta) + \tan(\theta)
    \qquad
    du = \big(\sec(\theta)\tan(\theta) + \sec^2(\theta)\big) d\theta
  \)}
  \[
    \onslide<+->{F = \int \frac{du}{u}}
    \onslide<+->{  = \ln\abs{u}+C}
    \onslide<+->{  = \ln \abs{\sec(\theta)+\tan(\theta)}+ C}
  \]
\end{frame}

%------------------------------------------------------------------------------%
\begin{frame}{Exemplo \theexemplo}
  \[
    \onslide<+->{x = a\sec(\theta)}
    \onslide<+->{\quad\Rightarrow\quad
    \sec(\theta)
      = \frac{\text{hipotenusa}}{\text{cateto adjacente}}
      = \frac{x}{a}}
  \]

  \addgraphic{6}{10,3}{exemplo_triangulo_3}

  \onslide<+->{\[
    \tan(\theta)
    = \frac{\text{cateto oposto}}{\text{cateto adjacente}}
    = \frac{\sqrt{x^2-a^2}}{a}
  \]}

  \begin{align*}
    \onslide<+->{F & = \int \frac{dx}{\sqrt{x^2-a^2}}                     \\[2mm]}
    \onslide<+->{  & = \ln \abs{\sec(\theta)+ \tan(\theta)} + C           \\[2mm]}
    \onslide<+->{  & = \ln \abs{\frac{x}{a}+ \frac{\sqrt{x^2-a^2}}{a}}+ C}
  \end{align*}
\end{frame}

%------------------------------------------------------------------------------%
\section{Lista Mínima}

%------------------------------------------------------------------------------%
\begin{frame}{Lista Mínima}
  Estudar a Seção 4.4 da Apostila

  Exercícios:

  \vfill
  Atenção: A prova é baseada no livro, não nas apresentações
\end{frame}

%------------------------------------------------------------------------------%
\end{document}
%------------------------------------------------------------------------------%